\chapter{Residual Boundary of the Arithmetic Branch}
\label{v4-arith:w10-final:chapter}

After the Koszul, chiral-homology, and arithmetic Chern--Simons
constructions have reached their natural boundary, the arithmetic
branch has named obstructions. One of them is equivalent to the
Riemann Hypothesis. The others are comparison, descent, regularity, and
small-prime hypotheses. The realization formalism of
Chapter~\ref{v4-ch:realization-formalism} is orthogonal to these
arithmetic conditions.

\section{Condensed-Hilbert Descent and Archimedean Tate Positivity}
\label{v4-arith:w10-final:w10-a}

The complete Weil form is defined on the compact convolution algebra
$\mathscr A$. A geometric Hilbert-space construction and a comparison
with this arithmetic form are separate data.

\begin{theorem}[Weil positivity and an additional descent]
\label{v4-arith:w10-final:thm:cs-descent-atp-rh}
For $\mathscr A=C_c^\infty(\mathbb R)$, positivity
$W(f*f^*)\ge0$ for every $f\in\mathscr A$ is equivalent to RH.
The same positivity together with a condensed-Hilbert descent therefore
implies RH. The descent is not needed for this analytic equivalence.
\end{theorem}
\begin{proof}
The explicit signed comparison \eqref{v4-arith:sc:full-comparison}
identifies the complete arithmetic form, with its two polar moments.
Weil's criterion proves the equivalence. Adding a further hypothesis
preserves the implication to RH.
\end{proof}

\begin{remark}[continuous and arithmetic purity]
\label{v4-arith:w10-final:rem:w10-a-shape}
Continuous scaling satisfies $|F_\infty|=e^{-1/2}1$ without RH.
The arithmetic zero operator satisfies $|F_Z|=e^{-1/2}1$ exactly
under RH, by Theorem~\ref{v4-arith:w9-r2:thm:purity-ATP-RH}.
A descent to the continuous carrier does not supply an intertwiner
with $F_Z$. That comparison must be specified and proved separately.
\end{remark}

\section{The RH-Equivalent Obstruction}
\label{v4-arith:w10-final:tier1}

The archimedean Tate positivity statement
(\Cref{v4-arith:w5-a:conj:archimedean-tate-positivity}) asserts:
\begin{equation}
\label{v4-arith:w10-final:eq:atp}
W(f*f^*)\ge0\qquad(f\in\mathscr A=C_c^\infty(\mathbb R)).
\end{equation}
Here $f^*(t)=\overline{f(-t)}$ and multiplication in $\mathscr A$
is additive convolution. Equivalently, on
$\mathscr P=\mathcal F_+\mathscr A$, the inequality is
$\mathcal W(hh^\sharp)\ge0$, where
$\mathcal W(h)=W(\mathcal F_+^{-1}h)$ and
$h^\sharp(z)=\overline{h(\bar z)}$.
Fubini gives $\mathcal F_+(f*f^*)=h h^\sharp$, so this is an
exact change of test algebra. The functional includes its two polar
moments, the signed gamma density, and the weights
$\Lambda(n)/\sqrt n$. equation~\eqref{v4-arith:sc:full-comparison} constructs
its signed pairing. Weil's criterion makes positivity for every test
equivalent to RH. It is denoted A-TP throughout.

\begin{proposition}[A-TP is an RH-equivalent obstruction]
\label{v4-arith:w10-final:prop:tier1}
\ClaimStatusProvedHere. A-TP
\eqref{v4-arith:w10-final:eq:atp} is an RH-equivalent obstruction in
the arithmetic branch.  No other residual in this chapter is used as an
RH-equivalent statement.
\end{proposition}

\begin{proof}
The Fourier map is an algebra isomorphism from convolution on
$\mathscr A$ to multiplication on $\mathscr P$ and respects the
stated involutions. Thus the displayed inequality is exactly Weil's
compact-test criterion. The additional comparison inputs do not enter
that equivalence.
\end{proof}

\section{Comparison and Descent Conjectures}
\label{v4-arith:w10-final:tier2}

\begin{definition}[comparison and descent inputs]
\label{v4-arith:w10-final:def:prior-table}
\ClaimStatusDefinitional. The comparison and descent boundary contains
the following named inputs.
\begin{center}
\small
\begin{tabular}{p{3.6cm}p{7.8cm}p{1.4cm}}
\toprule
\textbf{Input} & \textbf{Carries} & \textbf{Form}\\
\midrule
Wild-EG enhancement &
F1.a.2.class-match;
\(\mathrm{L}_{\mathrm{ACD}}\) Eisenstein \(\mathrm{Ext}^{1}\) at
\(\{0,1\}\); BC-glue wild inertia &
\(\mathrm{open}\)\\
\(\Pi\)-flat pro-\'etale flatness &
BC-diamond-glue cross-prime prismatic glue &
\(\mathrm{open}\)\\
Arithmetic AG singular support &
global arithmetic Arinkin--Gaitsgory singular-support condition and
the centre of BC-diamond-glue &
\(\mathrm{open}\)\\
Clausen--Scholze condensed-Hilbert descent &
Quantum arithmetic Chern--Simons bulk extension &
\(\mathrm{open}\)\\
Consani--Marcolli adele-class duality &
Positselski \(G\)-row duality and the abelian Heisenberg base-change
line &
\(\mathrm{open}\)\\
Chenevier pseudocharacter trace comparison &
Pairwise sheaf residual in the Ch1 closure &
\(\mathrm{open}\)\\
BCglue local comparison datum &
K5 and the Jacquet--Shalika trace comparison in BC-diamond-glue &
\(\mathrm{open}\)\\
Ben-Zvi lift at infinity &
\(\mathrm{L}_{\mathrm{ACD}}\) archimedean compatibility and finite
uniformity &
\(\mathrm{open}\)\\
BZN wild trace data &
Large-prime semisimplicity and the wild \(\mathrm{Ext}\)-trace pairing &
\(\mathrm{open}\)\\
Symmetric-square weight limit &
Weil--Petersson weight-aspect positivity for holomorphic cusp forms &
\(\mathrm{open}\)\\
Arithmetic bar nilpotency &
The square-zero property of the arithmetic pre-differential
\(d^{\mathrm{Ar}}\) &
\(\mathrm{open}\)\\
\bottomrule
\end{tabular}
\end{center}
\end{definition}

\begin{theorem}[visible non-RH residuals]
\label{v4-arith:w10-final:thm:four-prior}
\ClaimStatusProvedHere. Each input in
\Cref{v4-arith:w10-final:def:prior-table} is a non-RH residual named
elsewhere in the arithmetic branch. The statement does not assert
exhaustion of all possible residuals.
\end{theorem}

\begin{proof}
Wild-EG (the Emerton--Gee wild-inertia enhancement) carries F1.a.2 by
\Cref{v4-arith:w10-b:thm:class-match}, the
\(\mathrm{L}_{\mathrm{ACD}}\) exceptional Eisenstein crossing by
\Cref{v4-arith:w10-d:thm:class-match}, and the wild-inertia
component of BC-diamond-glue by
\Cref{v4-arith:w10-e:cor:bcglue-final-updated}.
\(\Pi\)-flat (pro-\'etale \(\Pi\)-flatness for the prismatic site) is
the cross-prime prismatic flatness of
\Cref{v4-arith:bcglue-k5:lem:prismatic-flatness}, unchanged by
Chapter~\ref{v4-arith:w10-e:chapter}. Arithmetic AG (arithmetic
Arinkin--Gaitsgory singular support) is the global singular-support
conjecture \Cref{v4-arith:conj:F1-global-AG}, required at the centre
of BC-diamond-glue by \Cref{v4-arith:w10-e:cor:bcglue-final-updated}.
Clausen--Scholze condensed-Hilbert descent supplies the Hilbert, trace,
and Wilson clauses of the chosen quantum ACS bulk model by
\Cref{v4-arith:w10-a:thm:descent-equivalent}.
The other rows are the named conditional inputs constructed in
Chapters~\ref{v4-arith:w5-d:chapter},
\ref{v4-arith:w5-h:bcglue-core-and-F2c},
\ref{v4-arith:bcglue-k5},
\ref{v4-arith:lacd-sub}, \ref{v4-arith:bzn-rcompletion},
\ref{v4-arith:w6-b:chapter}, and
\ref{v4-arith:arakelov-bar}.  Each is non-RH in form and is not
implied by A-TP.
\end{proof}

\begin{remark}[Wild-EG concentration]
\label{v4-arith:w10-final:rem:wild-eg}
Wild-EG is the only source conjecture carrying three distinct
class-match conditions. Its proof would simultaneously prove F1.a.2,
the exceptional \(\mathrm{L}_{\mathrm{ACD}}\) Eisenstein
\(\mathrm{Ext}^{1}\) crossing at \(\{0,1\}\), and the wild-inertia
component of BC-diamond-glue.
\end{remark}

\section{The Generalised Ramanujan Conjecture for
\texorpdfstring{\(\mathrm{GL}_{4}\)}{GL\_4}}
\label{v4-arith:w10-final:sub-rh}

The Rankin--Selberg square of a Maass form on \(\mathrm{GL}_{2}/\mathbb{Q}\)
comes from an isobaric automorphic object on
\(\mathrm{GL}_{4}/\mathbb{Q}\). Full temperedness of that
\(\mathrm{GL}_{4}\) object at every finite place, equivalently GRC for
its cuspidal factors, closes the Maass sector of the
\(\mathcal{V}^{\mathrm{prim}}\) Hochschild-trace identity. This
conjecture is separated from A-TP in the proof: it concerns Satake
eigenvalues, not zero placement.

\begin{hypothesis}[GRC for \(\mathrm{GL}_{4}/\mathbb{Q}\) in the Selberg class]
\label{v4-arith:w10-final:hyp:GRC-GL4}
\ClaimStatusConditional. The isobaric automorphic representations of
\(\mathrm{GL}_{4}/\mathbb{Q}\) arising from the Rankin--Selberg square
of a Maass form have tempered cuspidal factors at every finite place.
\end{hypothesis}

\begin{proposition}[GRC-\(\mathrm{GL}_{4}\) closes the Maass sector]
\label{v4-arith:w10-final:prop:GRC-GL4}
\ClaimStatusProvedHereConditional. Conditional on
\Cref{v4-arith:w10-final:hyp:GRC-GL4}, full Ramanujan holds for the
Maass forms entering the \(\mathcal{V}^{\mathrm{prim}}\)
Hochschild-trace \(P_{3}\) identity of \Cref{v4-arith:w10-c:thm:main}.
This closes the Maass sector. The hypothesis is not equivalent to RH;
no implication between this hypothesis and A-TP is used.
\end{proposition}

\begin{proof}
Chapter~\ref{v4-arith:w10-c:chapter} proves that the
Arthur--Ramakrishnan--GRC hypothesis implies full Ramanujan for Maass forms
(\Cref{v4-arith:w10-c:thm:AR-implies-Ramanujan}) and reduces that
hypothesis to \Cref{v4-arith:w10-final:hyp:GRC-GL4}
(\Cref{v4-arith:w10-c:prop:sub-RH-strict}). The \(P_{3}\) identity
is the exact conditional statement of \Cref{v4-arith:w10-c:thm:main}.
A-TP concerns the zeros of \(\zeta\) through Weil positivity;
GRC-\(\mathrm{GL}_{4}\) concerns Satake parameters.  The argument uses
the two residuals as separate inputs.
\end{proof}

\section{\(F_{1}\text{.a.1}\) Small Residuals}
\label{v4-arith:w10-final:tier3}

Two small residuals occur in the \(F_{1}\text{.a.1}\) table of
Chapter~\ref{v4-arith:w10-f:chapter}. They are not a count of all
small residuals in the arithmetic branch.

\begin{proposition}[small residuals]
\label{v4-arith:w10-final:prop:small-residuals}
\ClaimStatusProvedHere. Besides the generator-comparison inputs and the
automorphy/patching hypotheses isolated in
\Cref{v4-arith:w10-f:lem:BLGG-descent} and
\Cref{v4-arith:w10-f:lem:Emerton-patching},
the \(F_{1}\text{.a.1}\) table has the following two small residual tails:
\begin{enumerate}
\item \textbf{Supersingular small-weight tail}
(\Cref{v4-arith:w10-f:conj:ss-small-weight}): Fontaine--Laffaille
representations with Hodge--Tate weights
\(k \in \{p-1, p\}\), requires a non-formal theorem.
\item \textbf{\(A_{5}\) non-solvable descent defect}
(\Cref{v4-arith:w10-f:conj:ns-defect}): local-global descent defect
in the icosahedral sector, requires a non-formal theorem.
\end{enumerate}
\end{proposition}

\begin{proof}
The two listed residual tails are read from
\Cref{v4-arith:w10-f:prop:F1a1-table-w10f}. The Fontaine--Laffaille
supersingular core is reduced to the generator comparison of
\Cref{v4-arith:w10-f:thm:f1a1-ss-FL}. The odd-determinant icosahedral
core is reduced to the conditional theorem of
\Cref{v4-arith:w10-f:thm:A5-odd-closure}. The two tails listed above are
additional obstructions in that decomposition.
\end{proof}

\section{Separation from the Realization Formalism}
\label{v4-arith:w10-final:tier4}

The obstructions above are internal to the arithmetic branch. The
realization formalism of \Cref{v4-ch:realization-formalism} concerns
independent verification data for Vols I--III; the arithmetic branch
does not supply those data.

\begin{proposition}[realization formalism is orthogonal to the arithmetic branch]
\label{v4-arith:w10-final:prop:hziv-separate}
\ClaimStatusProvedHere. The arithmetic branch does not construct
independent comparison data for the source Vols I--III
ProvedHere claims. Its obstructions concern arithmetic comparison,
descent, temperedness, and small-prime regularity.
\end{proposition}

\begin{proof}
The preceding sections use only the arithmetic conditions A-TP,
comparison and descent, GRC for \(\mathrm{GL}_{4}\), and the two
\(F_{1}\text{.a.1}\) residuals. None of these data is a
independent comparison pair for a source theorem of Vols I--III.
\end{proof}

\section{Residual Boundary}
\label{v4-arith:w10-final:residual-boundary}

\begin{theorem}[arithmetic branch obstruction decomposition]
\label{v4-arith:w10-final:thm:near-closure}
\ClaimStatusProvedHere. The arithmetic branch has, at least, the
following obstructions:
\begin{enumerate}
\item one statement equivalent to RH: archimedean Tate positivity;
\item the comparison and descent inputs of
\Cref{v4-arith:w10-final:def:prior-table};
\item one Ramanujan-type hypothesis separated from A-TP: GRC for
\(\mathrm{GL}_{4}/\mathbb{Q}\);
\item the two \(F_{1}\text{.a.1}\) small residuals of
\Cref{v4-arith:w10-final:prop:small-residuals}.
\end{enumerate}
The Vol~IV realization formalism for Vols I--III is orthogonal to the
arithmetic branch.
\end{theorem}

\begin{proof}
Item (1) is \Cref{v4-arith:w10-final:prop:tier1}. Item (2) is
\Cref{v4-arith:w10-final:thm:four-prior}. Item (3) is
\Cref{v4-arith:w10-final:prop:GRC-GL4}. Item (4) is
\Cref{v4-arith:w10-final:prop:small-residuals}. The separation from
the realization formalism is \Cref{v4-arith:w10-final:prop:hziv-separate}.
No implication from the table to RH is asserted except through A-TP.
\end{proof}

\begin{remark}[domain of the boundary]
\label{v4-arith:w10-final:rem:not-rh}
The theorem isolates the displayed RH-equivalent content into A-TP and
places the visible non-RH obstructions in named inputs.  It does not
prove RH and it does not claim that the arithmetic comparison is
complete.
\end{remark}
